\documentclass[multi=tikzpicture,border=9pt]{standalone}
\usepackage[UTF8]{ctex}
\usepackage{amsmath,amssymb}
\usepackage{tikz}
\usetikzlibrary{arrows.meta,calc,patterns,decorations.pathmorphing}
\definecolor{main}{HTML}{17365D}
\definecolor{accent}{HTML}{AC3931}
\tikzset{every picture/.style={font=\small,line width=.75pt},
>={Stealth[length=2mm]},arr/.style={->},box/.style={draw=main,rounded corners=2pt,
fill=main!4,text width=2.7cm,minimum height=.9cm,align=center},
note/.style={align=center,text width=8cm}}
\begin{document}
\begin{tikzpicture}
\fill[main!10] (0,0) circle(1.8);\fill[white] (0,0) circle(1.3);\draw (0,0) circle(1.8);\draw (0,0) circle(1.3);\node at (0,.35){内腔 $P_2$};\node at (3.1,.8){外界 $P_0$};\node at (-2.7,.4){中间层 $P_1$};\draw[arr] (-2,.4)--(-1.5,.4);\draw[<->] (0,-.2)--(1.3,-.2) node[midway,below]{$a$};\draw[<->] (0,2.3)--(1.8,2.3) node[midway,above]{$b$};\draw[dashed] (1.8,0)--(1.8,2.3);\node at (-.9,1){$-\lambda$};\node at (1.7,1.25){$+\lambda$};\node[note] at (0,-2.7){横截面示意：长度足够大，薄膜无额外径向张力，间隙相对大小为示意。};
\end{tikzpicture}
\begin{tikzpicture}
\draw[arr] (0,0)--(6,0) node[right]{$r$};\draw[arr] (0,0)--(0,3.8) node[above]{$|E_r|$};\draw[main,thick] (0,0)--(1.5,0);\draw[main,thick] (1.50000,3.00000) -- (1.52000,2.97030) -- (1.54000,2.94118) -- (1.56000,2.91262) -- (1.58000,2.88462) -- (1.60000,2.85714) -- (1.62000,2.83019) -- (1.64000,2.80374) -- (1.66000,2.77778) -- (1.68000,2.75229) -- (1.70000,2.72727) -- (1.72000,2.70270) -- (1.74000,2.67857) -- (1.76000,2.65487) -- (1.78000,2.63158) -- (1.80000,2.60870) -- (1.82000,2.58621) -- (1.84000,2.56410) -- (1.86000,2.54237) -- (1.88000,2.52101) -- (1.90000,2.50000) -- (1.92000,2.47934) -- (1.94000,2.45902) -- (1.96000,2.43902) -- (1.98000,2.41935) -- (2.00000,2.40000) -- (2.02000,2.38095) -- (2.04000,2.36220) -- (2.06000,2.34375) -- (2.08000,2.32558) -- (2.10000,2.30769) -- (2.12000,2.29008) -- (2.14000,2.27273) -- (2.16000,2.25564) -- (2.18000,2.23881) -- (2.20000,2.22222) -- (2.22000,2.20588) -- (2.24000,2.18978) -- (2.26000,2.17391) -- (2.28000,2.15827) -- (2.30000,2.14286) -- (2.32000,2.12766) -- (2.34000,2.11268) -- (2.36000,2.09790) -- (2.38000,2.08333) -- (2.40000,2.06897) -- (2.42000,2.05479) -- (2.44000,2.04082) -- (2.46000,2.02703) -- (2.48000,2.01342) -- (2.50000,2.00000) -- (2.52000,1.98675) -- (2.54000,1.97368) -- (2.56000,1.96078) -- (2.58000,1.94805) -- (2.60000,1.93548) -- (2.62000,1.92308) -- (2.64000,1.91083) -- (2.66000,1.89873) -- (2.68000,1.88679) -- (2.70000,1.87500) -- (2.72000,1.86335) -- (2.74000,1.85185) -- (2.76000,1.84049) -- (2.78000,1.82927) -- (2.80000,1.81818) -- (2.82000,1.80723) -- (2.84000,1.79641) -- (2.86000,1.78571) -- (2.88000,1.77515) -- (2.90000,1.76471) -- (2.92000,1.75439) -- (2.94000,1.74419) -- (2.96000,1.73410) -- (2.98000,1.72414) -- (3.00000,1.71429) -- (3.02000,1.70455) -- (3.04000,1.69492) -- (3.06000,1.68539) -- (3.08000,1.67598) -- (3.10000,1.66667) -- (3.12000,1.65746) -- (3.14000,1.64835) -- (3.16000,1.63934) -- (3.18000,1.63043) -- (3.20000,1.62162) -- (3.22000,1.61290) -- (3.24000,1.60428) -- (3.26000,1.59574) -- (3.28000,1.58730) -- (3.30000,1.57895) -- (3.32000,1.57068) -- (3.34000,1.56250) -- (3.36000,1.55440) -- (3.38000,1.54639) -- (3.40000,1.53846) -- (3.42000,1.53061) -- (3.44000,1.52284) -- (3.46000,1.51515) -- (3.48000,1.50754) -- (3.50000,1.50000);\draw[main,thick] (3.5,0)--(5.6,0);\draw[dashed] (1.5,0)--(1.5,3);\draw[dashed] (3.5,0)--(3.5,1.5);\node[below] at (1.5,0){$a$};\node[below] at (3.5,0){$b$};\node[above] at (3.4,3){$E_r=-\lambda/(2\pi\varepsilon_0r)$};\node[note] at (2.6,-1.25){电场仅在两膜间存在，方向向内；这里画电场大小，未表示符号。};
\end{tikzpicture}
\begin{tikzpicture}
\draw[thick] (0,-1)--(0,1.5);\draw[thick] (5,-1)--(5,1.5);\draw[arr,accent] (0,.6)--(1.3,.6) node[right]{$p_{{\rm e},a}$};\draw[arr,accent] (5,.6)--(3.7,.6) node[left]{$p_{{\rm e},b}$};\node[above] at (0,1.5){内膜};\node[above] at (5,1.5){外膜};\node at (2.5,-.4){$p_{{\rm e},a}>p_{{\rm e},b}$};\node[note] at (2.5,-2){用两侧平均场：$p_{\rm e}=\varepsilon_0E^2/2$。内膜向外、外膜向内。};
\end{tikzpicture}
\begin{tikzpicture}
\node[box] (b0) at (0,0) {外膜：\\$P_1-P_0-p_{{\rm e},b}=0$};
\node[box] (b1) at (5,0) {$P_1=P_0+p_{{\rm e},b}$};
\node[box] (b2) at (5,-2) {内膜：\\$P_2-P_1+p_{{\rm e},a}=0$};
\node[box] (b3) at (0,-2) {$P_2=P_0+p_{{\rm e},b}-p_{{\rm e},a}$};
\draw[arr] (b0)--(b1);\draw[arr] (b1)--(b2);\draw[arr] (b2)--(b3);\node[note] at (2.5,-3.35) {内膜静电压力较大，故 $P_2<P_0$。薄间隙时差值是一阶曲率小量。};
\end{tikzpicture}
\begin{tikzpicture}
\node[box] (b0) at (0,0) {高斯定律求 $E_r$};
\node[box] (b1) at (5,0) {积分求 $\Delta\varphi$};
\node[box] (b2) at (5,-2) {平均场求膜压力};
\node[box] (b3) at (0,-2) {固定电荷能量导数核验};
\draw[arr] (b0)--(b1);\draw[arr] (b1)--(b2);\draw[arr] (b2)--(b3);\node[note] at (2.5,-3.35) {单位长度能量 $\lambda^2\ln(b/a)/(4\pi\varepsilon_0)$；固定电压变形须计入电源功。};
\end{tikzpicture}
\end{document}
